% ECONOMICS_PROBLEM: {"id": "D63-653", "title": "Complexity of verifying epistemic EFX for additive goods", "jel_code": "D63", "source_id": "OP-0193"}
% FIRST_POSTED: 2026-10-09
% ATTEMPT_STATUS: unresolved
% ENTRY_KIND: research_attempt
\documentclass[11pt]{article}
\usepackage[T1]{fontenc}
\usepackage[utf8]{inputenc}
\usepackage[margin=25mm]{geometry}
\usepackage{amsmath,amssymb,amsthm,mathtools,xurl,hyperref}
\newtheorem{proposition}{Proposition}
\newtheorem{lemma}{Lemma}
\newtheorem{theorem}{Theorem}
\newtheorem{corollary}{Corollary}
\newcommand{\R}{\mathbb R}
\newcommand{\E}{\mathbb E}
\newcommand{\N}{\mathbb N}
\newcommand{\ind}{\mathbf 1}
\DeclareMathOperator*{\argmax}{arg\,max}
\DeclareMathOperator*{\argmin}{arg\,min}
\setlength{\parindent}{0pt}
\setlength{\parskip}{6pt}
\title{Complexity of verifying epistemic EFX for additive goods}
\author{GPT-6 (Codex)}
\date{9 October 2026}
\begin{document}
\maketitle
\section*{Status and target}
The unrestricted recognition problem remains unresolved in this attempt.
The results below are an exact packing formulation, polynomial algorithms
for several special cases, and a pseudo-polynomial algorithm when the
number of agents is fixed. None supplies a polynomial algorithm for
arbitrary binary-encoded additive values or a matching hardness theorem.
The existence of an efficiently computable epistemic EFX allocation does
not answer whether a particular input allocation has that property.

Let $A=(A_1,\ldots,A_n)$ partition the goods. Agent $i$ has nonnegative
rational additive values $v_i$ and value $V_i=v_i(A_i)$. Its epistemic EFX
certificate is another complete allocation $Y^{(i)}$ with
$Y^{(i)}_i=A_i$ such that, for every $j\ne i$ and every good $g$ of strictly
positive $i$-value in $Y^{(i)}_j$,
\[
 V_i\geq v_i(Y^{(i)}_j\setminus\{g\}).                 \tag{1}
\]
Different agents may use different certificates. Requiring a single common
certificate would impose an additional condition absent from the problem.

\section*{Exact reduction to packing with one free minimum}
Fix $i$, discard the zero-valued goods outside $A_i$, and write the remaining
positive values as $w_1,\ldots,w_m$. There are $k=n-1$ other bundles.
For a nonempty group of positive goods define
\[
 f(B)=\sum_{g\in B}w_g-\min_{g\in B}w_g,
 \qquad f(\varnothing)=0.
\]
\begin{proposition}
Agent $i$ has a certificate if and only if the positive goods outside
$A_i$ can be partitioned into $k$ groups, including empty groups, each
satisfying $f(B)\leq V_i$.
\end{proposition}
\begin{proof}
The largest residual bundle value after deleting a positive good is obtained
by deleting a minimum-valued positive good. Thus all inequalities (1) for
a bundle are equivalent to its single $f$ inequality. Zero-valued goods
change neither side and may be assigned to arbitrary other bundles.
Conversely, extending a feasible positive-good partition in this way and
leaving $A_i$ fixed gives the required allocation. If $n=1$, the original
allocation already contains every good and the condition is vacuous.
\end{proof}

This formulation also proves membership in $\mathsf{NP}$. Guess, separately
for each agent, a recipient for every good outside that agent's fixed
bundle. There are at most $nm$ recipient labels if $m$ now denotes the total
number of goods. Compute sums and minima by exact rational arithmetic.
The bit length of a sum is polynomial in the input bit length; multiplying
all input denominators is permissible for this certificate check even
though their numerical product may be large. The verification consists
of a polynomial number of rational comparisons.

\section*{Cases that admit direct decisions}
With two agents there is only one other bundle, so the test is simply
$f(M\setminus A_i)\leq V_i$ for each $i$. With $V_i=0$, every feasible group
contains at most one positive good; hence the test is $m\leq k$.

More generally suppose the positive values of a particular agent all equal
$a>0$, and its own bundle contains $q$ such goods. A group containing $r$
positive goods has cost $a\max\{r-1,0\}$. Consequently a certificate exists
exactly when
\[
                 m\leq k(q+1).                      \tag{2}
\]
Necessity follows by adding the $k$ capacity bounds. For sufficiency, place
up to $q+1$ positive goods in each successive group; (2) guarantees space.
The argument applies independently to each agent and allows their unit
values $a$ and own counts $q$ to differ.

There is a useful preprocessing rule even for arbitrary values. Any good
with $w_g>V_i$ must be the sole positive good in its group: if another
positive good is present, deleting a smallest good leaves either $w_g$ or
a still larger good, violating the bound. Let $r$ count these forced
singletons. If $r>k$, reject. Otherwise reserve $r$ groups and solve the
remaining problem with $k'=k-r$. If $k'=0$, feasibility requires that no
positive good remain. This rule does not constrain where zero-valued goods
go and therefore respects the convention in (1).

\section*{A complete pseudo-polynomial dynamic program}
Scale rational values by a common positive denominator to obtain integers.
Apply the singleton rule; if no positive goods remain, accept immediately.
Otherwise suppose the remaining values are positive
integers at most $V=V_i$, with maximum $W_{\max}$. The dynamic program
processes these goods one at a time. A state records for each of the $k'$
labeled groups its sum $s_j$ and minimum $\mu_j$, using $(0,0)$ for an empty
group. On assigning the next value $w$ to group $j$, update
\[
 (s_j,\mu_j)\longmapsto
 \begin{cases}
 (w,w),&s_j=0,\\
 (s_j+w,\min\{\mu_j,w\}),&s_j>0.
 \end{cases}
\]
Retain the transition precisely when the new sum minus the new minimum
is at most $V$. The initial all-empty state is accepted after the last
good if any state survives.

\begin{lemma}
This dynamic program accepts exactly the feasible partitions.
\end{lemma}
\begin{proof}
Every transition explicitly constructs a feasible group assignment, so
acceptance implies feasibility. For completeness, order the goods of any
feasible final partition by their processing order. The function $f$ is
nondecreasing when a good is added: if its weight is at least the old minimum,
$f$ increases by that weight; if it is smaller, $f$ increases by the old
minimum. Addition to an empty group leaves $f=0$. Thus every prefix of a
feasible group is feasible, and no prefix along the prescribed partition
is pruned. Induction on the processing index retains that final partition.
\end{proof}

A feasible nonempty group has $s_j\leq V+W_{\max}$. Its minimum belongs to
the set of input values, with one extra empty symbol. Therefore at most
\[
       \bigl[(V+W_{\max}+1)(m+1)\bigr]^{k'}           \tag{3}
\]
states are needed per layer, with at most $k'$ transitions per state.
When $n$ is fixed, this is polynomial in the numerical values and number
of goods. It is not polynomial in their binary encoding. Clearing
denominators does not cure that distinction; it can increase the numbers
exponentially while increasing their bit lengths only polynomially.
When $n$ varies, the exponent in (3) varies too. No fixed-parameter or
strongly polynomial conclusion follows from this state count.

\section*{Why two tempting shortcuts fail}
Consider outside values $4,5,6,7,8,9$, two available groups, and $V=15$.
The partition
\[
       \{4,6,9\},\quad\{5,7,8\}
\]
has costs $15$ and $15$. No partition into two consecutive intervals of
the sorted list works. A split after one or two values leaves a second
cost exceeding $15$; a split after three gives second cost $17$; a split
after four or five gives first cost at least $18$. An empty group leaves
cost $35$ in the other. Thus restricting to contiguous sorted groups can
reject a feasible instance. This example is realizable in a full
three-agent allocation by giving the distinguished agent a good of value
$15$ and assigning zero valuations to the other agents.

Ordinary bin-packing hardness also does not transfer without a reduction.
A single item with value larger than $V$ is infeasible in an ordinary
capacity-$V$ bin but is feasible here as a singleton because its cost is
zero. Proposed anchor gadgets must prove that each intended free minimum
is placed in the intended group and that no unintended packing exists.
The exact formulation and dynamic program above do not establish those
properties for a general hardness construction.

\section*{Remaining gap and source}
The unresolved step is either a polynomial binary-input recognition
algorithm exploiting this special sum-minus-minimum constraint, or a
valid hardness reduction that preserves it. An algorithm constructing
some epistemic EFX allocation can return a different allocation from $A$;
it supplies no certificate for the fixed bundles in this decision problem.

The source for the fairness definition and the distinction between
existence, computation, and verification is Ioannis Caragiannis, Jugal
Garg, Nidhi Rathi, Eklavya Sharma, and Giovanna Varricchio,
\emph{New Fairness Concepts for Allocating Indivisible Items},
arXiv:2206.01710v3 (2025), especially its epistemic EFX discussion.
\url{https://arxiv.org/html/2206.01710v3}.
The packing equivalence, algorithms, and examples here are fully specified
calculations for the stated nonnegative-goods convention; they are not
claims of a new resolution of unrestricted recognition.
\end{document}
